\documentclass[12pt]{article}

\title{The axiom of well-ordered choice fails in New Foundations}

\author{Lavinia Randall Holmes}

\usepackage{amssymb}

\begin{document}


\maketitle

These are notes on Elliott Glazer's disproof of the Axiom of Well-Ordered Choice (hereinafter AC$_{\tt wo}$) in New Foundations.   

The Specker tree of a cardinal $\kappa$ is defined as the collection of all cardinals $\lambda$ such that $\exp^n(\lambda)= \kappa$ for some $n$.

Writing down things I know.  That Specker trees are well-founded (due to Forster;  this justifies calling these structures trees) relies on an identity of the form $\aleph(\kappa) \leq \exp^i(\kappa)$.  If we are high enough in the tree, $i$ can be taken to be 2.  If we are two levels up from the bottom of the Specker tree, we have a countable subset.  Thus if we ($\kappa$) are four levels up from the bottom of the Specker tree we have $\kappa^2=\kappa$ (because we can construct Quine pairs).  Thus well orderings of $X \in \kappa$ can be represented in ${\cal P}(X)$, and order types of well-orderings of subsets of $X$ in ${\cal P}^2(X)$, so $\aleph(\kappa) \leq {\tt exp}^2(\kappa)$.  Since $\aleph({\tt exp}^2(\kappa)) \not\leq {\tt exp}^2(\kappa)$, we have $\aleph(\kappa)<\aleph({\tt exp}^2(\kappa))$, and from this it follows that Specker trees must be well-founded, without any appeal to choice.  Consider a minimal
$\aleph(\kappa)$ for a cardinal $\kappa$ at least four levels up in the Specker tree.  This $\kappa$ cannot itself be an image
under ${\tt exp}^2$ (say ${\tt exp}^2(\lambda)$) of a $\lambda$ at least four levels up in the Specker tree, because we would have $\aleph(\lambda) < \aleph({\tt exp}^2(\lambda)) = \alpha(\kappa)$ contrary to choice of $\aleph(\kappa)$.  So such
a $\kappa$ cannot be more than five steps from the bottom of the Specker tree along any path, and the tree is well-founded.

We assume AC$_{\tt wo}$, the proposition that any well-ordered partition has a choice set.

Under this assumption we demonstrate a bound on $\aleph({\tt exp}(\kappa))$ (for $\kappa$ high enough in a Specker tree that Suppose $X \in \kappa$ and $\alpha$ is the order type of a well-ordering of ${\cal P}(X)$ with domain $Q$  We can from the well ordering given
get a well-ordering of $Q \times Q$.  With each pair $(A,B)$ in the domain of this well ordering for which $A \Delta B$ is nonempty, associate the collection of ordered pairs $(x,y)$ where $x \in A \setminus B$ and $y \in B\setminus A$, when both of these sets are nonempty, or $x=y \in A \Delta B$ when one of the two sets is empty.  Use well-ordered choice to associate a single such pair $(x,y)$ with each such pair $A\times B$, a relation $R$ on whose pairs we have a well ordering.
Then replace each set in the domain of the well-ordering taken from ${\cal P}(X)$ with its intersection with the field of the relation $R$.  This replacement is injective, because every difference between two sets in the domain of the well-ordering is witnessed by at least one element of the field of $R$.  Further, the cut-down sets are well-orderable, and can be placed in an injective correspondence with subsets of an ordinal less than $\alpha^2\cdot 2$.
So $\aleph({\exp}(\kappa)) \leq \aleph(\exp(\aleph(\kappa)))$, which gives us the desired bound. 

Now to do sophisticated bounding.  What we really want is a uniform bound on the lengths of all paths in the Specker tree of a cardinal $\lambda$.

$\aleph({\tt exp}^n(\kappa))$ should be boundable.   It is bounded above by $\aleph((\exp\circ \aleph)^{n}(\kappa)))$ by repeated application of the bounding identity, that is $(\aleph \circ \exp)^n(\aleph(\kappa))$.   $(\aleph \circ \exp)$ is a monotone increasing operation on alephs.  Moreover, it sends alephs to larger alephs.  $\aleph((\exp\circ \aleph)^{n}(\kappa)))$ itself is boundable.  It is less than or equal to
${\tt exp}^{3n+2}(\kappa)$ and so strictly less than $\aleph({\tt exp}^{3n+2}(\kappa))$.

What this indicates is that if $\lambda$ is a cardinal and $\kappa$ is in its Specker tree with $\aleph(\kappa)$ minimal and $\lambda = \exp^n(\kappa)$, that $\aleph(\lambda)$ is less than 
$\aleph((\exp\circ \aleph)^{n}(\kappa)))$ which is in turn less than $\aleph({\tt exp}^{3n+2}(\kappa'))$ for any $\kappa'$ whatsoever in the Specker tree, so the tree is flat.

The above naively assumes that all exponentials exist.  But such an assumption is not needed.  If we assume that the rank of the Specker tree is infinite, we can specialize to
$\kappa'$ for which the indicated exponential exists, because it is far enough down in the tree.  Strictly speaking there is a similar concern about existence of the given bound on $\aleph(\lambda)$.
But choose a cardinal at level $3n+2$ or more down from the top of the tree.  Its Hartogs aleph will have the bound indicated computable (at $n$ steps upward), which means the bound for $\aleph(\kappa)$ itself was computable.

$\aleph(T^n(|V|) = \aleph(\exp(T^{n-1}(|V|)) \leq \aleph(\exp(\aleph(T^{n-1}(|V|)))$.

This is enough to establish a contradiction.  $\aleph(T^n(|V|))$ is strictly less than $\aleph(\aleph(\exp(\aleph(T^{n-1}(V)))))$.  We have a function $f(\beta)$ which commutes with T ($f(\beta)=\aleph(\aleph(\exp(\beta))))$) is monotone and $f(\beta)>\beta$, and an ordinal $\alpha$ ($\aleph(T^{n-1}(V))$) such that $f(T(\alpha)) > \alpha$.

Now consider the smallest ordinal $\gamma$ such that for some $m$, $f^m(\gamma)\geq \alpha$.
$T\gamma$ is the smallest ordinal $\delta$ such that for some $m$, $f^m(\delta)\geq T(\alpha)$, but this is clearly
$\gamma$ itself.

Now consider the smallest $m$ such that $f^m(\gamma)\geq\alpha$.  Clearly the smallest $n$ such that
$f^n(\gamma)\geq T(\alpha)$ is $Tm$ and the smallest $n$ such that  $f^n(\gamma)\geq T^2(\alpha)$ is $T^2m$.  But then $m$ exceeds either $Tm$ or $T^2m$ (in a bizarre but notionally possible case) by one and this is impossible.

So AC$_{\tt wo}$ is false in NF.

My technique above adapts to prove ESI.

Suppose that $X$ is an infinite  set and $C$ is a well-ordering of subsets of $X$.  Then $C \times C$ can be well-ordered.

Go through $C\times C$, which has order type $\alpha$ constructing a function from $1+\alpha=\alpha$ to subsets of $X$.  First, put the collection of all elements of $X$ which do not belong to any symmetric difference of elements of $C$. Then In connection with each pair of sets $(A,B)$, put in $A \setminus B$ minus the union of all sets previously considered.  The nonempty sets in the range of this function make up a well-ordered partition of a subset of $X$.  The sets in C are injectively associated with unions of sets in this partition.  So each set in $C$ is associated with a subset of $\aleph^*(X)$ [the set of ordinal indices of included elements of the partition], and we have
shown that each well-ordering of subsets of $C$ is associated with a well-ordering of subsets of the power set of the order type of a partition of $X$, so $\aleph(\exp(X)) \leq \aleph(\exp(\aleph(\exp(\aleph^*(X)))))$.

As soon as PP is assumed, we get a bounding inequality from which we can prove that Specker trees are flat much as above and from which we can prove a contradiction in NF just as above, using a different $f$ easily read from the different form of ESI.




\end{document}